\documentclass[border=10pt]{standalone}
\usepackage{fontspec}
\setmainfont{Noto Serif}
\newfontfamily\bengalifont{Noto Serif Bengali}[Script=Bengali]
\usepackage[Latin,Bengali]{ucharclasses}
\setTransitionsFor{Bengali}{\begingroup\bengalifont}{\endgroup}
\usepackage{newunicodechar}
\newunicodechar{।}{{\bengalifont\char"0964}}
\newunicodechar{॥}{{\bengalifont\char"0965}}

\usepackage[table]{xcolor}
\usepackage{booktabs,tabularx,array}
\definecolor{bapilrule}{HTML}{1F2937}
\definecolor{bapilstripe}{HTML}{F3F4F6}
\renewcommand{\arraystretch}{1.28}
\setlength{\tabcolsep}{5pt}

\begin{document}
\begin{minipage}{17cm}
{\large\bfseries সারণি ৬৮।}\ {\large সংশোধনের ইতিহাস: অক্টোবর ২০২৬-এর নিরীক্ষার পর করা পরিবর্তন}\par\vspace{7pt}
\small
\rowcolors{2}{bapilstripe}{white}
\begin{tabularx}{\linewidth}{@{}>{\hsize=0.667\hsize\raggedright\arraybackslash}X>{\hsize=1.500\hsize\raggedright\arraybackslash}X>{\hsize=1.083\hsize\raggedright\arraybackslash}X>{\hsize=0.750\hsize\raggedright\arraybackslash}X@{}}
\arrayrulecolor{bapilrule}\toprule
\rowcolor{white}\textbf{কোথায়} & \textbf{পরিবর্তন} & \textbf{কারণ} & \textbf{প্রমাণ বা হিসাব} \\
\midrule
সূচনা-অংশ & অবস্থা-সারণি; দায়স্বীকার ও এক নজরে অংশ সংশোধিত; এক নজরে অংশে বাজেট যুক্ত; সংখ্যা হালনাগাদ & যা বিদ্যমান আর যা প্রস্তাবিত তা আলাদা করা & রিপোজিটরি যাচাই, ২২ অক্টোবর ২০২৬; সংযুক্তি ত \\
অধ্যায় ১ & আইনের সময়রেখা শর্তযুক্ত; পদ্ধতির ধাপ ৬ (নিরীক্ষা); বাস্তবায়নকে গবেষণা-প্রোটোটাইপ হিসেবে বর্ণনা & আইনগত ও সফটওয়্যারের অবস্থা অতিরঞ্জন এড়ানো & [84, 85, 86, 113] \\
অধ্যায় ২; সারণি ৪ & EGDI ও GCI-এর প্রাথমিক উৎস; ২০২৩ সালের এআই প্রস্তুতি-র‍্যাংকের তথ্যসূত্র সংশোধন; আইনগত সারাংশ পুনর্লিখন & নির্ভুলতা ও উৎসের মান & [80, 119, 120] \\
অধ্যায় ৩ & এআই হাবকে বাস্তবায়নাধীন প্রকল্প নয়, স্বাক্ষরিত অনুদান-চুক্তিসহ একটি প্রস্তাব হিসেবে বর্ণনা & উৎসে অনুমোদনের অপেক্ষায় থাকা প্রস্তাবের কথা & [35, 117] \\
অধ্যায় ৫; সংযুক্তি ঙ থেকে ছ & সংস্থাগুলোর আইনগত অবস্থান; ERRC সদস্যদের স্বাধীন বাছাই; কেবল যুক্তিসংগত কারণে অপসারণ; আইনের প্রাধান্য (ধারা ১৩); আইনি পর্যালোচনার পর সনদ নিশ্চিতকরণ & স্বাধীনতা ও আইনগত নির্ভুলতা & সংযুক্তি ধ \\
অধ্যায় ৮, ১০ & প্রধান ডেটা কর্মকর্তা, NRDEX ও সোর্স কোড ভান্ডার যে দলিলে বর্ণিত তার নামে উল্লিখিত; ডেটা দেশে রাখার সংশোধনী; কঠোরতর ব্যাখ্যার নিয়ম; GPU-সক্ষমতা বাজেটের সঙ্গে সংগতিপূর্ণ & সঠিক উৎস-নির্দেশ; অভ্যন্তরীণ সংগতি & [84, 113, 115]; সারণি ৪৯ \\
অধ্যায় ৯; সারণি ২৮, ৩০ & বেঞ্চমার্কের গঠন, বিভাজন ও পরিসংখ্যানগত যুক্তি; সিদ্ধান্ত-নিয়ম; মাপকাঠির সংজ্ঞা পরিমার্জন; সম্ভাব্যতা-প্রাক্কলনসহ পাঁচ ধাপের মডেল-ধাপক্রম & লক্ষ্য পরিমাপযোগ্য করা; নতুন মডেল প্রয়োজন বলে ধরে না নেওয়া & উইলসন ব্যবধান; [121] \\
অধ্যায় ১১; SECURITY.md & পরিপক্বতার পরিভাষা সংজ্ঞায়িত; CI নির্ধারিত বলে বর্ণিত; নিরাপত্তা ও নির্ভরতা-চর্চা & প্রোটোটাইপ ও উৎপাদন-ব্যবস্থা আলাদা করা & রিপোজিটরি \\
অধ্যায় ১২, ১৩ & P1-এর উত্তরণ প্রটোকলের সঙ্গে যুক্ত; অভিযোগ ও আপিল; লেবেলকে BAPIL-এর নিয়ম হিসেবে উপস্থাপন & পরিমাপযোগ্যতা; জবাবদিহি; সঠিক উৎস-নির্দেশ & [89]; সংযুক্তি দ \\
অধ্যায় ১৪; সারণি ৪১ থেকে ৪৩ & অধ্যাদেশ ও আইনের সময়ক্রম; প্রণয়ন ও কার্যকারিতা আলাদা; আইনি পর্যালোচনার পরিধি; নতুন অনুচ্ছেদ ১৪.৭ & আইনগত বক্তব্য মাধ্যমিক উৎসনির্ভর & [84 থেকে 89, 106, 113 থেকে 116] \\
অধ্যায় ১৫; সারণি ৪৪, ৪৮ & সমঝোতা স্মারক, আয়োজক-চুক্তি ও অর্থায়ন দিয়ে মাস ১ সংজ্ঞায়িত; MS1 ও প্রথম ৯০ দিন সংগতিপূর্ণ; ক্যালেন্ডার-সতর্কতা & বাস্তবায়নের বাস্তবতা & সংযুক্তি ন \\
অধ্যায় ১৭; সারণি ৪৯, ৫০ & বাজেট ১৮.৪৯ থেকে ২০.৪৩ কোটি টাকায় সংশোধিত: মূল্যবৃদ্ধি, ধাপে ধাপে নিয়োগ, পৃথক আপৎকালীন সংস্থান, মূল্যায়নের উপাদান; পরিস্থিতি ও সংবেদনশীলতা; আসন্নীকরণের ব্যাখ্যা & আর্থিক স্বচ্ছতা & budget\_model.py; [118]; সারণি ৬১ \\
অধ্যায় ১৮, ১৯; সারণি ১৬, ১৮, ২৫, ৫৪ & KPI লক্ষ্য হিসেবে চিহ্নিত; সাশ্রয়কৃত সময় স্ব-প্রতিবেদন নয়, পরিমাপকৃত; প্রশিক্ষণের সংখ্যা বাজেটের সঙ্গে মেলানো; OCR নমুনা; ঝুঁকি পর্যবেক্ষণের সংযোগ & পরিকল্পনা, বাজেট ও পরিমাপের সংগতি & সংযুক্তি দ; সারণি ৪৯ \\
অধ্যায় ২১; সারণি ৫৭ & উপসংহার পুনর্লিখিত; সিদ্ধান্ত D4 সংশোধিত; কোনো সিদ্ধান্ত সংবিধিবদ্ধ সংস্থা সৃষ্টি করে না বলে টীকা & অতিরঞ্জন; স্বাধীনতা & সংযুক্তি ধ \\
সংযুক্তি ত থেকে ফ & প্রমাণগত অবস্থা; বাজেটের মিলকরণ ও পরিস্থিতি; মূল্যায়ন প্রটোকল; শাসন-অবস্থা; পূর্বশর্ত ও কাজের প্যাকেজ; ঝুঁকি পর্যবেক্ষণ; এই সংশোধনের ইতিহাস ও অমীমাংসিত বিষয় & পর্যালোচনার নির্দেশনা & প্রতিটি সংযুক্তিতে উল্লিখিত \\
তথ্যসূত্র & তথ্যসূত্র ১১৩ থেকে ১২১ যুক্ত & নতুন প্রমাণ & তালিকা অনুযায়ী \\
বাংলা সংস্করণ & সব সংশোধন অনূদিত; তালিকা, উদ্ধৃতি, পৃষ্ঠা-নম্বর ও চার্টের লেবেলে বাংলা সংখ্যা & দ্বিভাষিক সংগতি; সম্পাদকীয় মান & অংশভিত্তিক সংখ্যার স্বয়ংক্রিয় তুলনা \\
\bottomrule
\end{tabularx}
\par\vspace{5pt}{\footnotesize\itshape পূর্ণ সমস্যা-নিবন্ধন, দাবি-যাচাই ম্যাট্রিক্স ও স্কোর রিপোজিটরির পৃথক নিরীক্ষা ও সংশোধন প্রতিবেদনে (blueprint/\allowbreak{}review/\allowbreak{}) আছে।}
\par\vspace{3pt}{\scriptsize\color{gray} উৎস: BAPIL কারিগরি ব্লুপ্রিন্ট, সংস্করণ ১.০ (২০২৬)। CC BY 4.0।}
\end{minipage}
\end{document}
